\section{システム境界と脅威}\label{sec:overview}

\subsection{システム境界と主体}\label{sec:overview:layout}

Constitutional ARI-RQGM が統治するのは、自律的な研究実行系を取り巻く
制度である。研究実行系は、型の定まった成果物、測定値、査読記録、原稿
候補を保存点へ出力する。統治機構はそれらを観測し、どの指示文、構成
要素、効用規則を稼働させるかを決める。ホスト計算機の基本ソフトウェア、
モデルの重み、固定プログラムを直接変更できる人間は統治の対象外とする。

主体は四群に分かれる。研究担当は提案、評価、弁護、執筆を行う。司法
担当は証拠収集、訴追、弁護、裁定を行う。メタ担当は指示文や効用規則の
候補を作り、過去事例から再試験用の事例を選ぶ。固定担当は憲法規則を
検査し、台帳の遷移を適用し、影響を受けた記録を修復し、論文の最終主張
を検証する。各担当がハッシュ値付きの記録を作れるのは、あらかじめ宣言
された役割の範囲内だけである。現実装は作成者を認証する電子署名や鍵
管理を提供しない。助言を出す権限と状態を変更する権限は分離する。

\subsection{脅威モデル}\label{sec:overview:pipeline}

中心的な脅威は、制度が自分自身を再帰的に最適化することである。更新
可能な役割が研究の信頼性より自己保存を優先し、点数を水増しし、評価
指標を攻略し、証拠を省き、入力文に埋め込まれた命令へ従う可能性がある。
提案役と評価役の共謀、現職より広い権限を求めるメタ担当、自分への訴追
を自ら裁く裁定役、退役後も派生した点数を通じて選択へ影響する役割も
想定する。

偶発的な故障も対象に含める。台帳更新の途中で処理が停止する、モデルが
形式に合わない出力を返す、比較用の基準資料がない、指示文候補が出力
形式に違反する、再採点に必要な評価軸が保存されていない、といった場合
である。このとき研究実行は継続できても、制度変更は安全側に倒して拒否
しなければならない。

悪意ある管理者、固定プログラムの任意改変、保存点と局所固定値を一緒に
過去へ戻す攻撃、基礎的なハッシュ関数の破綻、提供者によるモデル重みの
秘密変更は対象外である。外部 API、実験装置、本番データへの不可逆操作
もカーネルの統治対象外であり、別の固定承認門を要する。憲法プログラム
の変更は人間による改正であって、稼働中の仕組みだけでは防げない。

\subsection{権限と検証境界}\label{sec:overview:verification}

権限は一方向に流れる。モデルを使う役割は、提案、攻撃、弁護、裁定案、
後継候補を生成できる。しかし、それらを直接台帳の状態へ変換することは
できない。決定論的カーネルが、記録、実行役、要求された遷移、実行時期、
権限の包含関係、来歴、監査履歴の連続性を検査する。台帳の唯一の書き手
は遷移処理であり、その出力も確定前に再検査される。

研究実行と制度変更では、故障時の方針を分ける。研究経路上の統治処理が
失敗した場合も、候補生成と局所計算は続ける。これは隔離実行環境内で
外部副作用がない場合の可用性方針であり、現実装はその条件自体をカーネル
で強制しない。状態を変更する統治処理が失敗した場合は、証拠と適法性を
確認できない限り変更を拒否し、現職制度を維持する。

最終的な原稿検証も固定境界に属する。原稿が前向きに宣言した構造化数値
主張を実行証拠から再計算し、必要な証拠が揃っているかを調べる。一般の
自然言語から主張を意味的に抽出する処理ではない。執筆役と査読役は原稿
候補を選び改稿できるが、検証器と判定結果は変更できない。

\subsection{状態、来歴、再現}\label{sec:overview:checkpoint}

保存点は、一つの制度史を識別する単位である。追記専用の出来事列を正本
とし、台帳、期の状態、影響除去の集計は、出来事列から再構築できる派生
情報とする。統治対象の記録には、作成された期、稼働中だった構成要素、
指示文と効用規則のハッシュ値を付ける。監査記録は直前の記録を参照する
ハッシュ連鎖をなし、内容改変と後続を残した途中削除を検出する。残存する
局所固定値より短い履歴も検出できるが、履歴と固定値を同じ過去状態へ戻す
末尾切り詰めは、外部の最新先頭値がないため検出できない。

期の更新は、準備、旧期の終了、新期の開始、確定という四段階で記録する。
再開時には未完の末尾を無視し、境界処理を決定論的にやり直す。創設登録
も一度だけ確定し、確定済みなら記録から再構築するだけで再実行しない。
論理的な影響除去ではデータを削除せず、古いことと選択対象外であること
を追記し、なぜ退役の影響を受けたかという証拠を残す。

期の識別値は二つに分けてから合成する。政策指紋には、稼働制度と、
しきい値・予算を含む解決済みの統治設定を結ぶ。実行指紋には、宣言された
モデル、復号設定、道具、環境、データの固定版を結ぶ。提供者や環境の厳密な版
が得られない場合は未解決と明記する。前者は制度間の比較に使い、後者は
同じ政策が異なる宣言済み実行基盤で動いたことを見分ける。不透明なモデル
サービス自体の再現性を保証するものではない。

\subsection{設計原則}\label{sec:overview:principles}

設計判断は次の五原則に要約できる。

\begin{enumerate}
  \item \textbf{固定法と進化する制度。}憲法、遷移規則、最終検証経路は、
        稼働中に更新できる成果物の集合へ入れない。
  \item \textbf{期の間に制度を固定。}期が開いている間は稼働中の版を
        変えない。重大違反の T16 隔離は直ちに境界を形成し、新しい指紋値
        の期を開くため、その前後には異なる制度識別子が付く。
  \item \textbf{助言と執行を分離。}提案役、裁定役、メタ担当は記録を
        出力するだけとし、台帳はカーネル検証下の遷移処理だけが書く。
  \item \textbf{履歴を残し、影響だけを断つ。}退役しても証拠は消さない。
        ただし、その役割に依存する子孫記録は再評価するか選択対象外に
        してから、将来の選択へ反映する。
  \item \textbf{証拠に見合う主張。}構造上の保証と、動作確認と、研究
        品質の改善を別々に扱う。機構が動いたことを品質向上の証拠とは
        みなさない。
\end{enumerate}
